\chapter{Formulations of Classical Mechanics}
\label{fm}

Classical mechanics is usually presented as the most settled part of
physics: three laws, a variational principle, and a semester of
worked examples. By the time a student reaches quantum theory or
general relativity, mechanics has already receded into the
background as a body of technique rather than a live subject. This
impression is not wrong, exactly, but it is incomplete. 

\section {(In)equivalence of the Formulations} 
\label{fm:nf}

The Newtonian, Lagrangian, Hamiltonian, and Hamilton--Jacobi formulations of
classical mechanics are routinely presented as equivalent packagings of a
single physical theory, related by Hamilton's principle, the Legendre
transform, and Jacobi's theorem. This claim is correct, but only under
conditions that are rarely stated explicitly and that fail even for
holonomic systems considered above.

Here we show a number  of
 distinct ways in which the standard equivalence breaks down: 
\begin{itemize}
   \item the gap
between Hamilton's principle (stationarity of the action) and the popularly
stated ``principle of \emph{least} action'' (minimality), which fails beyond a
conjugate time fixed by the curvature of the potential and which can be made
arbitrarily short;
   \item the consequent failure of the boundary-value problem built
into Hamilton's principle itself to be well-posed---even for $H=p^2/2m+V(q)$, a
sufficiently steep $V$ makes the two-endpoint problem admit no solution, or
infinitely many, in contrast to the always-well-posed Newtonian initial-value
problem;
   \item the existence of singular Lagrangians, for which the Legendre
transform is not invertible and Dirac's theory of constraints is required;
   \item the existence of Newtonian force laws that admit no
Lagrangian description at all, governed by the Helmholtz conditions;
   \item the
failure of the Legendre transform to be a global diffeomorphism even for
pointwise-regular Lagrangians, absent the stronger condition of
hyperregularity;
   \item the mechanistically identical failure of the Hamilton--Jacobi
equation's generating function to stay single-valued beyond a caustic;
\end{itemize}

Let us display shortly the four standard formulations.

\paragraph{Newtonian mechanics}
For a system of $N$ particles with positions $\bm{r}_i \in \mathbb{R}^3$,
Newton's second law reads 
\[m_i \ddot{\bm{r}}_i = \bm{F}_i,\]
where $\bm{F}_i$ is the total force on particle $i$. Constraints are handled by
splitting the force into an applied part and a constraint (reaction) part,
\[\bm{F}_i = \bm{F}_i^{\mathrm{app}} + \bm{F}_i^{\mathrm{c}},\]
and postulating---this is d'Alembert's principle---that the constraint forces 
do no virtual work: 
\[\sum_i \bm{F}_i^{\mathrm{c}}\cdot \delta \bm{r}_i = 0\]
for every virtual displacement $\delta\bm{r}_i$ compatible with the
constraints at a frozen instant of time. It is d'Alembert's principle, not
Newton's second law by itself, that supplies the bridge to the variational
formulations below.

\paragraph{Lagrangian mechanics}
Let $Q$ be the configuration manifold and $TQ$ its tangent bundle, with induced
coordinates $(q^i,\dot q^i)$. A Lagrangian is a function 
\[L : TQ \times \mathbb{R} \to \mathbb{R},\]
typically (but not necessarily) of the form $L = T- V$, 
kinetic minus potential energy. Hamilton's principle asserts that
physical trajectories $q(t)$ extremize the action
\begin{equation}
S[q(\cdot)] = \int_{t_0}^{t_1} L\big(q(t),\dot q(t), t\big)\, dt
\end{equation}
among curves with fixed endpoints $q(t_0), q(t_1)$. Setting the first variation
to zero yields the Euler--Lagrange equations
\begin{equation}
\label{EL}
\frac{d}{dt}\frac{\partial L}{\partial\dot q^i} - \frac{\partial L}{\partial q^i} = 0 ,
    \qquad i = 1,\dots,n=\dim Q .
\end{equation}
For a system subject to \emph{ideal, holonomic} constraints, this is fully
equivalent to Newton's second law plus d'Alembert's principle, once $Q$ is
taken to be the constraint surface itself and the $q^i$ are independent
generalized coordinates on it.

\paragraph{Hamiltonian mechanics}
Define the momenta 
\[p_i = \partial L/\partial \dot q^i.\]
This is the \emph{fiber derivative} 
\[\mathbb{F}L: TQ \to T^*Q \quad (\bm q,\dot {\bm q})\mapsto (\bm q, \bm p).\]
Where $\mathbb{F}L$ can be inverted for $\dot {\bm q}$ in terms of $(\bm q,\bm p)$,\
one defines the Hamiltonian by the Legendre transform,
\begin{equation}
    H(\bm q,\bm p,t) = p_i \dot q^i - 
    L(\bm q,\dot {\bm q}, t)\Big|_{\dot q = \dot q(q,p,t)} ,
\end{equation}
and Hamilton's canonical equations \footnote{ These equations also can 
be obtained from variational principle discussed later}
\begin{equation}
\label{Ham}
\dot q^i = \frac{\partial H}{\partial p_i}, \qquad 
    \dot p_i = -\frac{\partial H}{\partial q^i}
\end{equation}
reproduce~\eqref{EL} on the nose. The cotangent bundle $T^*Q$ carries a
canonical symplectic form $\omega = dp_i \wedge dq^i$, giving Hamiltonian
mechanics its distinctive geometry: canonical transformations, Poisson
brackets, Liouville's theorem, and so on.

The precise statement licensing the textbook claim is the following.
\begin{proposition}[Regular equivalence]
\label{regular}
Let $$L: TQ\times\mathbb{R}\to\mathbb{R}$$
    be \emph{regular}: the fiber Hessian
    $$W_{ij} = \partial^2 L/\partial \dot q^i \partial \dot q^j$$
    is nondegenerate at every point of $TQ$. Then $\mathbb{F}L$ is a local
    diffeomorphism, $H$ is well defined locally, and solutions of~\eqref{EL}
    correspond, via 
    $$p = \mathbb{F}L(q,\dot q),$$ 
    bijectively and locally to
    solutions of~\eqref{Ham}. If in addition $\mathbb{F}L$ is a \emph{global}
    diffeomorphism $TQ \to T^*Q$---in which case $L$ is called
    \emph{hyperregular}---this correspondence extends to a global bijection
    between the full solution spaces of the two theories.
\end{proposition}

\paragraph{Hamilton--Jacobi theory}
A fourth formulation, historically inseparable from Hamilton's own work,
replaces the $2n$ first-order equations~\eqref{Ham} by a single first-order
partial differential equation for a scalar function. Define \emph{Hamilton's
principal function} $S(q,t)$ as the action evaluated along the true trajectory
of the system that starts at some fixed reference event $(q_0,t_0)$ and ends at
$(q,t)$. Hamilton showed that, viewed as a function of its endpoint, $S$
satisfies the \emph{Hamilton--Jacobi equation}
\begin{equation}
\label{HJ}
\frac{\partial S}{\partial t}+H\Big(q,\frac{\partial S}{\partial q},t\Big)=0.
\end{equation}
Conversely, Jacobi showed how to run this construction backwards and recover
the canonical equations~\eqref{Ham} directly from~\eqref{HJ}, without
integrating an ordinary differential equation, provided one can find a
\emph{complete integral}: an $n$-parameter family of solutions $S(q,t;\alpha)$,
$\alpha=(\alpha_1,\dots,\alpha_n)$, of~\eqref{HJ} whose mixed Hessian
$\partial^2 S/\partial q^i\partial \alpha_j$ is everywhere nondegenerate on the
\emph{domain of interest}.

\begin{proposition}[Jacobi's theorem]
\label{jacobi-thm}
Let $S(q,t;\alpha)$ be a complete integral of~\eqref{HJ} on an open set
    $U\subset Q\times\mathbb{R}$, with $\det(\partial^2 S/\partial
    q^i\partial\alpha_j)\neq 0$ throughout $U$. Then the equations
\begin{equation}
\frac{\partial S}{\partial\alpha_j} = \beta_j \ (\text{constants}), 
    \qquad p_i = \frac{\partial S}{\partial q^i}
\end{equation}
implicitly define, for each choice of constants $(\alpha,\beta)$, a curve
    $(q(t),p(t))$ solving Hamilton's equations~\eqref{Ham} on $U$; every
    solution of~\eqref{Ham} passing through $U$ arises this way for a suitable
    choice of $(\alpha,\beta)$.
\end{proposition}

Proposition~\ref{jacobi-thm} is the formal engine behind solving
Hamiltonian systems by separation of variables, and it is a genuine equivalence
result: wherever a complete integral with nondegenerate mixed Hessian exists,
the Hamilton--Jacobi and Hamiltonian pictures determine exactly the same
trajectories. The qualification \emph{'on the domain of interest'} is doing real
work, and  cannot in general be
dropped, and that its failure is tied to the same mechanism that undermines the
``least action'' reading of Hamilton's principle.

Everything that follows  concerns ways in which the
hypotheses of Proposition~\ref{regular} or
Proposition~\ref{jacobi-thm} can fail to hold, or ways in which their
being satisfied does not settle every reasonable sense in which one might want
two physical theories to be ``the same.'' 

The equivalence between Hamilton's principle and the Euler--Lagrange
equations \eqref{EL} is exact and unconditional: $\delta S=0$ for all
admissible variations if and only if \eqref{EL} holds along the trajectory.
This is a first-order, local statement about the action, and none of
standard formulations depends on anything more. But the popular name for
Hamilton's principle---the ``principle of \emph{least} action''---asserts
something stronger: that the true trajectory does not merely make $S$
stationary but \emph{minimizes} it among nearby competitors with the same
endpoints. This stronger claim concerns the second variation, not the first,
and it is generally false, even for the most elementary holonomic,
unconstrained, regular Lagrangian system one could write down.

For a one-dimensional natural Lagrangian 
$$L(x,\dot x) = \tfrac12 m\dot x^2 - V(x),$$
consider a perturbation $x(t)+\varepsilon\eta(t)$ of a solution $x(t)$
of \eqref{EL}, with $\eta(t_0)=\eta(t_1)=0$. The second variation of the action
is
\begin{equation}
\delta^2 S[\eta] = \int_{t_0}^{t_1} \Big[ m\dot\eta(t)^2 - V''(x(t))\,\eta(t)^2 \Big]\, dt ,
\end{equation}
and the trajectory is a genuine local minimum of $S$ if and only if
$\delta^2S[\eta]>0$ for every admissible $\eta\not\equiv0$. Whether this holds
is governed by the \emph{Jacobi accessory equation}, obtained by
linearizing~\eqref{EL} about $x(t)$,
\begin{equation}
\label{jacobi-eq}
m\ddot\eta + V''(x(t))\,\eta = 0 ,
\end{equation}
together with $\eta(t_0)=0$. A time $t_c>t_0$ is a \emph{conjugate point} (or
\emph{kinetic focus}) to $t_0$ along the trajectory if the nontrivial solution
of~\eqref{jacobi-eq} with $\eta(t_0)=0$ vanishes again at $t_c$.

\begin{proposition}[Jacobi's necessary condition]
\label{jacobi-necessary}
The trajectory $x(t)$ furnishes a local minimum of the action on $[t_0,t_1]$
    only if $t_1$ occurs before the first conjugate point to $t_0$. If $t_1$
    lies beyond the first conjugate point, $S$ is a saddle point at $x(t)$:
    nearby trial trajectories exist with strictly smaller action, even though
    $x(t)$ exactly satisfies~\eqref{EL}.
\end{proposition}

\begin{example}[Harmonic oscillator]
\label{sho}
Take $V(x)=\tfrac12 kx^2$, so $V''\equiv k$ and \eqref{jacobi-eq} becomes
    $\ddot\eta+\omega^2\eta=0$ with $\omega=\sqrt{k/m}$. The solution with
    $\eta(t_0)=0$ is $\eta(t)=\sin\!\big(\omega(t-t_0)\big)$, which next
    vanishes at $t_c=t_0+\pi/\omega$---exactly half a period. By
    Proposition~\ref{jacobi-necessary}, a segment of oscillator trajectory
    genuinely minimizes the action if its duration is less than $\pi/\omega$,
    and is merely a stationary saddle point if its duration exceeds
    $\pi/\omega$.
\end{example}

Example~\ref{sho} answers, quantitatively, the question of how short is ``short
enough.'' The conjugate-point time $\pi/\omega=\pi\sqrt{m/k}$ is set entirely
by the curvature $V''=k$ of the potential and shrinks without bound as $k$
grows: no fixed time interval is safe for every system, because a sufficiently
steeply curved potential will already have produced a conjugate point within
it. More generally, conjugate points arise only where $V''>0$ along the
trajectory; for potentials with $V''(x)\leq 0$ for all $x$ (a free particle,
uniform gravity, an inverted quadratic barrier) no conjugate point ever forms,
and the action is minimized for arbitrarily long times.
Whether ``the action is minimized'' is therefore not a property of Hamilton's
principle in general---it is a property of the specific potential, and
specifically of how sharply it curves.

It is worth being precise about what this does and does not show. Hamilton's
principle, in its correct form ($\delta S=0$), remains exactly equivalent to
Newton's second law throughout; nothing here contradicts
Proposition~\ref{regular}. What fails is the stronger, popularly stated
identification of Hamilton's principle with a principle of \emph{least}
action---and it fails already in the cleanest setting available anywhere in
this paper: a single unconstrained particle, regular and hyperregular
Lagrangian, smooth potential. That is precisely why this failure mode belongs
at the head of this subsection, ahead of the more exotic constructions that
follow.

The question is asked whether a trajectory already known to
satisfy~\eqref{EL} minimizes the action. A logically prior question is whether
such a trajectory exists and is unique at all, given the data Hamilton's
principle is actually stated in terms of: not an initial position and velocity,
but two fixed \emph{endpoints}, $q(t_0)=q_A$ and $q(t_1)=q_B$---exactly the
boundary conditions built into the action functional. 
This is a fundamentally different mathematical
problem from the Newtonian initial-value problem. Given $(q_0,\dot q_0)$ at a
single instant, standard ODE theory (Picard--Lindelöf) guarantees a unique
solution of~\eqref{EL}, at least locally in time, for any reasonably smooth
force law: the ``most boring'' Hamiltonian, $H=p^2/2m+V(q)$, is always
well-posed as an initial-value problem. Fixing $q(t_0)$ and $q(t_1)$ instead
offers no such guarantee. The boundary-value problem can have no solution, a
unique solution, or infinitely many, and which of these holds is governed by
exactly the conjugate points.

Fix $t_0$ and $q_A=q(t_0)$, and let $x(t,v_0)$ denote the unique initial-value
solution with $\dot x(t_0)=v_0$. The \emph{shooting map} $v_0\mapsto
x(t_1,v_0)$ determines which endpoints are reachable in time $T=t_1-t_0$ and
how many initial velocities reach a given one. Its derivative $\partial
x(t,v_0)/\partial v_0$ satisfies~\eqref{EL} linearized about the
trajectory---for $H=p^2/2m+V(q)$ this is exactly the Jacobi accessory
equation~\eqref{jacobi-eq}, with initial data $\eta(t_0)=0$, $\dot\eta(t_0)=1$.
Consequently
\begin{equation}
\frac{\partial x(t_1,v_0)}{\partial v_0} = 0
    \quad\Longleftrightarrow\quad t_1 \text{ is a conjugate point to } t_0 ,
\end{equation}
so the shooting map is a local diffeomorphism precisely up to the first
conjugate time, and can fold---losing injectivity, surjectivity, or both---at
and beyond it.

\begin{example}[Harmonic oscillator boundary-value problem]
\label{sho-bvp}
Continuing Example~\ref{sho}, set $t_0=0$, $q_A=0$. The initial-value solution
    is $x(t,v_0)=(v_0/\omega)\sin(\omega t)$, so the shooting map to time $T$
    is $v_0\mapsto (v_0/\omega)\sin(\omega T)$, and
\begin{equation}
v_0 = \frac{\omega\, q_B}{\sin(\omega T)}
\end{equation}
solves $x(0)=0$, $x(T)=q_B$ whenever $\sin(\omega T)\neq0$: a unique solution,
    as expected for $T$ short of the first conjugate time $\pi/\omega$. But at
    $T=n\pi/\omega$ for a positive integer $n$---an integer multiple of the
    conjugate time---$\sin(\omega T)=0$ identically, and the shooting map
    collapses to $v_0\mapsto 0$ regardless of $v_0$: for $q_B\neq0$ \emph{no}
    solution exists, while for $q_B=0$ \emph{every} $v_0\in\mathbb{R}$ solves
    the boundary-value problem, a one-parameter family of distinct true
    trajectories connecting the same two events.
\end{example}

Example~\ref{sho-bvp} shows non-existence and non-uniqueness sitting on
opposite sides of the same coincidence, both triggered by the same
conjugate-time condition, and both occurring for the most elementary confining
potential there is. The critical travel times $n\pi/\omega$ are evenly spaced
only because the harmonic oscillator's own equation of motion is already
linear---so it coincides with its own Jacobi equation, and conjugate points
recur exactly periodically. For a general potential, the accessory
equation~\eqref{jacobi-eq} has a time-varying coefficient $V''(x(t))/m$ and the
conjugate times need not be evenly spaced, but the qualitative picture
survives, and can be more dramatic still: so quartic
oscillator, $V(x)=\tfrac14 Cx^4$, in which a whole family of distinct true
trajectories leaving a common starting event recross one another at a common
later event, producing multiple genuinely different---not merely differently
scaled---solutions of the same boundary-value problem.

 Since $\omega=\sqrt{k/m}$ for the harmonic case,
the critical travel times $n\pi/\omega = n\pi\sqrt{m/k}$ shrink as the
curvature $k$ grows, so for \emph{any} fixed travel time $T$, a steep enough
potential well places $T$ at or past a degenerate configuration. ``The most
boring system there is,'' $H=p^2/2m+V(q)$, therefore fails to have a well-posed
boundary-value problem in exactly the regime---strongly confining, rapidly
growing $V$---that most textbook potentials are drawn from.

This multiplicity is not a curiosity confined to classical mechanics: it is
exactly what the Feynman path-integral formulation of quantum mechanics must
sum over. The semiclassical propagator between $q_A$ and $q_B$ is, in general,
a sum of contributions from \emph{every} classical trajectory solving the
boundary-value problem, each carrying a phase determined by its action together
with a further phase set by the \emph{Maslov index}---the number of conjugate
points the trajectory has crossed. The Maslov index is, in
this sense, a bookkeeping device  it is needed because the classical
boundary-value problem it
quantizes is not, in general, uniquely solvable.

A Lagrangian is \emph{singular} (or degenerate) when its fiber Hessian $W_{ij}
= \partial^2 L/\partial \dot q^i \partial \dot q^j$ is degenerate somewhere on
$TQ$. In that case $\mathbb{F}L$ is not a local diffeomorphism there: the
momenta $p_i = \partial L/\partial \dot q^i$ are not independent functions of
the velocities, and one cannot solve for all of $\dot q$ in terms of $(q,p)$.
Concretely, degeneracy of $W_{ij}$ means the map $\mathbb{F}L$ satisfies one or
more relations
\begin{equation}
\phi_a(q,p) = 0, \qquad a = 1,\dots,k,
\end{equation}
purely among the momenta and coordinates---\emph{primary constraints}, in
Dirac's terminology---that hold identically on the image of $\mathbb{F}L$,
before any equations of motion are imposed.

The cleanest illustration is to note that Hamiltonian mechanics is itself,
formally, a maximally singular special case of Lagrangian mechanics. Treat
$(q,p)$ as $2n$ independent coordinates on an auxiliary configuration space and
consider the first-order ``Hamiltonian-form'' Lagrangian
\begin{equation}
L_1(q,p,\dot q,\dot p) = p_i \dot q^i - H(q,p).
\end{equation}
Its Euler--Lagrange equations reproduce~\eqref{Ham} exactly, so in one sense
Hamiltonian mechanics \emph{is} a special case of Lagrangian mechanics on a
bigger space. But the fiber Hessian of $L_1$ with respect to $(\dot q,\dot p)$
vanishes identically---$L_1$ is linear, not quadratic, in the velocities---so
$L_1$ is about as singular as a Lagrangian can be, and the naive
Legendre-transform machinery simply does not
apply to it. Recovering Hamilton's equations from $L_1$ requires running the
full constrained-Hamiltonian algorithm below, not the regular Legendre
transform.

The general treatment of singular Lagrangians is due to Dirac and Bergmann.
 One passes to the \emph{primary constraint surface}
$\Sigma_1 \subset T^*Q$ cut out by the $\phi_a$, and checks whether the primary
constraints are preserved under time evolution generated by a (now only
partially determined) Hamiltonian; if not, this consistency requirement
generates \emph{secondary constraints}, and the process is iterated until a
stable constraint surface is reached. Constraints are further classified as
\emph{first class} (those with vanishing Poisson brackets with all other
constraints, which typically signal gauge freedom---directions in phase space
that are not physically distinguishable) or \emph{second class} (which can be
eliminated using the Dirac bracket, an amended Poisson bracket adapted to the
constraint surface). The upshot is that for a
singular $L$, the physical phase space is not $T^*Q$ but a further-reduced
submanifold, of dimension strictly smaller than $2n$ in the presence of
first-class constraints, and the map from Lagrangian solutions to Hamiltonian
solutions is not the clean bijection of Proposition~\ref{regular} but a
considerably more elaborate correspondence, often carrying residual gauge
freedom on the Lagrangian side that has no counterpart in a naively
written-down Hamiltonian system.

Singular Lagrangians are not a pathology confined to toy examples: they are the
generic situation for gauge theories. Electromagnetism, Yang--Mills theory, and
general relativity all have singular Lagrangians, precisely because the local
gauge symmetry (arising from Noether's second theorem) forces degeneracy of the
fiber Hessian. For essentially all of fundamental field theory, then, the
``equivalence'' between the Lagrangian and Hamiltonian pictures is not the
elementary Legendre-transform story at all, but
the considerably more delicate Dirac--Bergmann correspondence.

\paragraph{The inverse problem: not every Newtonian system is Lagrangian}

Even in the unconstrained case, the passage from Newtonian to Lagrangian
mechanics is not guaranteed to exist at all. Given a system of second-order
ODEs $\ddot q^i = f^i(q,\dot q,t)$---an arbitrary Newtonian force law, divided
by mass---one can ask whether it arises as the Euler--Lagrange equations of
\emph{some} Lagrangian. This is the \emph{inverse problem of the calculus of
variations}. The answer is governed by the \emph{Helmholtz conditions}: the
linearized (Fr\'echet) differential operator associated with the equations of
motion must be self-adjoint. Douglas (1941) gave a complete classification
of when a Lagrangian exists in the two-degree-of-freedom case, building on
earlier work of Hirsch and Mayer; the general $n$-dimensional theory remains an
active area of research in the geometric calculus of variations.

The familiar case that fails the Helmholtz conditions is linear
velocity-dependent damping. The damped harmonic oscillator,
\begin{equation}
m\ddot x + \gamma \dot x + k x = 0 ,
\end{equation}
is not the Euler--Lagrange equation of any Lagrangian of the ordinary
autonomous, single-particle form $L(x,\dot x)$: the friction term breaks the
self-adjointness that the Helmholtz conditions require. This is not merely a
failure of imagination; it reflects the fact that a Hamilton's-principle
Lagrangian for a single degree of freedom automatically conserves a symplectic
volume (Liouville's theorem), while a damped oscillator's phase-space volume
visibly contracts as the system loses energy to friction. One can restore a
variational description, but only at a structural cost: by doubling the degrees
of freedom (the Bateman construction, pairing the damped oscillator with a
time-reversed, amplified partner), by using an explicitly time-dependent
Lagrangian such as the Caldirola--Kanai form $L = e^{\gamma t/m}\big(\tfrac12
m\dot x^2 - \tfrac12 kx^2\big)$, or by admitting non-standard
(non-quadratic-in-velocity) Lagrangians. Each of these routes reproduces the
correct equation of motion, but none of them is the ``obvious,''
structure-preserving Lagrangian one would have guessed, and none is unique.

The general moral is that Newtonian mechanics, understood as ``any second-order
ODE system one likes,'' is strictly more permissive than Lagrangian mechanics:
the class of Lagrangian-representable force laws is a proper,
Helmholtz-constrained subset of the class of Newtonian force laws, unless one
is willing to pay for a Lagrangian description with auxiliary variables or
explicit time dependence.

\paragraph{Global obstructions: regularity is not hyperregularity}
Proposition~\ref{regular} distinguishes \emph{regularity} (the fiber
Hessian is nondegenerate at each point) from \emph{hyperregularity} (the fiber
derivative $\mathbb{F}L$ is a diffeomorphism of the whole of $TQ$ onto the
whole of $T^*Q$). The distinction matters because regularity is a purely local,
pointwise condition, while the equivalence of the full \emph{solution spaces}
of the Lagrangian and Hamiltonian theories is a global statement about
$\mathbb{F}L$.

A Lagrangian can be regular everywhere and still fail to be hyperregular, in
either of two ways. First, $\mathbb{F}L$ can fail to be \emph{surjective}: if
$L$ grows sub-quadratically in $\dot q$ at large velocity (rather than the
quadratic-plus-lower-order growth of the usual kinetic-minus-potential form),
the momenta $p=\partial L/\partial\dot q$ can be bounded even as $\dot
q\to\infty$, so that some momentum values in $T^*Q$ are simply never reached,
and the corresponding Hamiltonian trajectories have no Lagrangian counterpart.
Second, $\mathbb{F}L$ can fail to be \emph{injective} globally even while being
a local diffeomorphism everywhere---the fiber derivative can ``fold back'' on
itself far from any particular point, so that two distinct velocities at the
same base point map to the same momentum. In either case, one is forced either
to restrict attention to a proper open subset of phase space on which
$\mathbb{F}L$ does behave well, or to accept that the Lagrangian and
Hamiltonian pictures describe genuinely different (if overlapping) collections
of physically possible histories. This is a comparatively mild failure mode,
 but it shows that even the
textbook regularity condition, usually checked (if at all) by verifying $\det
W_{ij}\neq 0$ at a single representative point, is not by itself sufficient for
the global equivalence the folklore claims.

\paragraph{The Hamilton--Jacobi correspondence: caustics}
Jacobi's theorem (Proposition~\ref{jacobi-thm}) is, like
Proposition~\ref{regular}, a genuine equivalence result---but, like
Proposition~\ref{regular} without hyperregularity, it is only a
\emph{local} one, and the nondegeneracy hypothesis on the mixed Hessian
$\partial^2S/\partial q\partial\alpha$ is not guaranteed to persist away from
the point where it is first checked.

Construct Hamilton's principal function concretely,
fix an initial event $(q_0,t_0)$, and for each
nearby $(q,t)$ let $S(q,t)$ be the action along the trajectory that solves
Hamilton's equations with $q(t_0)=q_0$ and $q(t)=q$. For $S$ to be a
well-defined, single-valued, smooth function of $(q,t)$, this connecting
trajectory must exist and be \emph{unique}. Existence is guaranteed for short
enough time by ordinary differential equation theory. Uniqueness is exactly
where the conjugate points re-enter: consider the
family of trajectories leaving $q_0$ at $t_0$ with varying initial velocity,
and ask where two members of this family, with infinitesimally different
initial velocity, next intersect. This is precisely the kinetic-focus
construction, and it recurs at the same time $t_c$
computed there. Beyond $t_c$, distinct nearby trajectories from $(q_0,t_0)$ can
reconverge near the same later point, so more than one candidate value competes
to be $S(q,t)$: the function becomes multivalued, its would-be smooth branches
meeting along an envelope---the \emph{caustic}---where the Jacobian of the map
from initial velocity to endpoint, and with it the mixed Hessian of
Proposition~\ref{jacobi-thm}, degenerates.

\begin{proposition}
\label{caustic}
Let $t_c$ be the first conjugate time to $t_0$ along a trajectory.
    Hamilton's principal function $S(q,t)$, built
    from trajectories issuing from $(q_0,t_0)$, is a smooth, single-valued
    solution of the Hamilton--Jacobi equation~\eqref{HJ} for $t<t_c$, and
    generically fails to remain single-valued for $t>t_c$, where the family of
    trajectories develops a caustic.
\end{proposition}

Nothing dynamically singular happens to the particle at a caustic: the actual
trajectories---the characteristics of~\eqref{HJ}---pass through it perfectly
smoothly, exactly as they pass through the conjugate point.
What breaks down is the specific representational
device of Hamilton--Jacobi theory: a single, smooth, globally valid generating
function from which the entire flow can be read off via $p=\partial S/\partial
q$. This is the same phenomenon, and the same underlying mathematics, that
produces caustics in geometric optics---where Hamilton first developed the
parallel theory---and it is the classical root of the phase corrections (the
Maslov index) that enter semiclassical, WKB-type quantization once a trajectory
has passed through one or more foci.

``The Hamiltonian and
Hamilton--Jacobi formulations are equivalent'' is true, without qualification,
only up to the first conjugate point; beyond it, recovering the dynamics from
Hamilton--Jacobi theory requires patching together multiple local branches of
$S$, or abandoning the single smooth generating function in favor of directly
integrating~\eqref{Ham}. This has nothing to do
with singular Lagrangians, nonholonomic constraints, or a failure of
hyperregularity: the same one-dimensional particle in a smooth potential well
illustrates this failure too, because
a conjugate point in the variational sense and a point on a caustic in the
Hamilton--Jacobi sense are the same event, described in two different
formulations of classical mechanics.

\paragraph{Non-uniqueness: the correspondence is not one-to-one in the other direction}
Even restricting entirely to the well-behaved case---holonomic constraints, $L$
regular and hyperregular---the map from ``the physics'' to ``a Lagrangian (or
Hamiltonian) representing it'' is many-to-one rather than one-to-one, which
complicates any claim that a given Lagrangian and a given Hamiltonian
formulation are simply \emph{the} representations of one theory. Adding a total
time derivative, $L \mapsto L + \frac{d}{dt}F(q,t)$ for arbitrary $F$, leaves
the Euler--Lagrange equations, and hence all physical predictions, unchanged;
more drastically, so-called $s$-equivalent Lagrangians can differ by more than
a total derivative and still generate the same force law, again governed by the
multiplier version of the Helmholtz analysis. On
the Hamiltonian side, canonical transformations relate infinitely many
distinct-looking triples $(q,p,H)$ that all describe the same underlying
dynamics. None of this contradicts Proposition~\ref{regular}---the
\emph{solution sets} still correspond bijectively---but it means the intuitive
picture of ``one Newtonian system, one Lagrangian, one Hamiltonian''
understates how much freedom is built into the correspondence in both
directions.







\section{Hertz's forceless mechanics }
\label{frm:hertz}

By the 1890s, several physicists --- Kirchhoff and Mach among them
--- had grown uneasy with the status of \emph{force} as a
fundamental concept of mechanics. Force is defined operationally
through $F=ma$, but that same equation is routinely used to
\emph{introduce} force, which makes the law read uncomfortably like
a definition dressed up as an empirical statement. In his
posthumously published \emph{Die Prinzipien der Mechanik in neuem
Zusammenhange dargestellt} (1894), Heinrich Hertz proposed to remove
the discomfort by removing force altogether: he set out to build the
whole of mechanics from just three primitive notions --- space,
time, and mass --- with no independent concept of force at
all~\cite{hertz}.

\subsection{The principle of the straightest path}

Consider $N$ mass points with positions $\mathbf x_i\in\mathbb R^3$
and masses $m_i$, and equip the $3N$-dimensional configuration space
$Q$ with the natural mass-weighted (kinetic-energy) metric,
\begin{equation}
  ds^2 \;=\; \sum_{i=1}^N m_i\, d\mathbf x_i\cdot d\mathbf x_i,
  \qquad
  T \;=\; \frac12\left(\frac{ds}{dt}\right)^{\!2}.
  \label{hertz-metric}
\end{equation}
Hertz's masses are never free in the naive sense: they are always
tied together by \emph{rigid, workless connections} (idealized
mechanical linkages that transmit force but do no work), which
restrict the accessible configurations to some submanifold
$\mathcal C\subset Q$. His fundamental postulate --- the
\emph{principle of the straightest path}
--- is then simply:

\begin{quote}
\emph{A system free of "force" moves along the straightest possible
path compatible with its constraints, at constant speed.}
\end{quote}

"Straightest possible" is to be read literally, in the geometric
sense: the representative point traces a \emph{geodesic} of the
metric~\eqref{hertz-metric} restricted to $\mathcal C$, i.e.\ a
curve whose acceleration (in the induced metric) is purely normal to
$\mathcal C$ and is entirely accounted for by the constraint
reaction. Formally, if the constraints are written $g_a(\mathbf
x)=0$, $a=1,\dots,k$, the equations of motion are
\begin{equation}
  m_i\,\ddot{\mathbf x}_i \;=\; \sum_a \lambda_a\,
  \frac{\partial g_a}{\partial \mathbf x_i},
  \label{hertz-eom}
\end{equation}
with the multipliers $\lambda_a$ fixed by demanding $g_a(\mathbf
x(t))\equiv 0$. There is nothing on the right-hand side of
\eqref{hertz-eom} except constraint reactions: no applied-force
term appears anywhere, by construction. Constant speed
($|ds/dt|=\mathrm{const}$, i.e.\ constant kinetic energy) follows
automatically, since ideal constraints do no work.

\subsection{Hidden masses, or where "force" goes}

Of course, ordinary bodies \emph{do} appear to exert forces on one
another --- gravity, contact forces, and so on all look like genuine
interactions between visible masses that are not rigidly connected
to anything. Hertz's resolution was to posit \emph{hidden masses}:
invisible mass points, rigidly linked to the visible ones (and, in
his more speculative moments, identified with a mechanical
aether), whose only role is to curve the constraint submanifold
$\mathcal C$ in just the right way. If the full configuration space
splits as $Q = Q_{\mathrm{vis}}\times Q_{\mathrm{hid}}$ and the
\emph{complete} system --- visible plus hidden --- moves along a
geodesic of $\mathcal C$, then projecting that geodesic motion down
onto $Q_{\mathrm{vis}}$ alone generically produces a curve that is
\emph{not} geodesic in the visible sector by itself: extra terms
appear in its equation of motion that are, from the point of view of
an observer with access only to $Q_{\mathrm{vis}}$, indistinguishable
from the effect of an applied force~\cite{hertz}.

This is a genuinely striking claim: on Hertz's account, force is not
absent from experience, only absent from the fundamental ontology.
Every observed force law would, in principle, be explained by some
specific (if permanently unobservable) arrangement of hidden mass
and rigid linkage. The same basic move --- explain an apparent force
as the shadow, on a submanifold, of pure geometric (geodesic) motion
in a larger space --- reappears twice in twentieth-century physics
in far more successful forms: in Kaluza--Klein theory, where the
electromagnetic force on a charged particle is recovered as geodesic
motion in a five-dimensional space, and, far more importantly, in
general relativity, where gravity itself --- which Newton had to
treat as a force acting at a distance --- becomes geodesic motion of
a free particle through curved four-dimensional spacetime, with no
"force" required at all. Hertz did not anticipate curved spacetime;
his geometry is Euclidean and the curvature is manufactured entirely
by hidden extra coordinates. But the structural resemblance to
Einstein's later, physically successful geometrization of gravity is
close enough that Hertz's book is still read today chiefly for that
reason, rather than as a working theory in its own
right.

\subsection{A second, more conservative geometrization}

It is worth noting that ordinary, force-\emph{having} mechanics can
also be written as a geodesic principle, without any hidden masses
at all, via Jacobi's form of the principle of least action. For a
conservative system with potential $V(\mathbf x)$ and fixed total
energy $E$, the true trajectories are exactly the geodesics of the
\emph{Jacobi metric}
\begin{equation*}
  ds_J^2 \;=\; 2\bigl(E-V(\mathbf x)\bigr)\, ds^2 ,
\end{equation*}
a conformal rescaling of the mass metric~\eqref{hertz-metric} by
the local kinetic energy. So there are, in the end, two quite
different ways to eliminate force from the description of a system
governed by a potential: enlarge the configuration space with hidden
coordinates and keep the metric flat (Hertz), or keep the visible
configuration space as it is and let the metric itself become curved
and energy-dependent (Jacobi). Both reproduce exactly the same
observable trajectories. Neither is more "true" than the other; they
are, in Hertz's own vocabulary, different \emph{images} of the same
phenomena. That a concept as apparently basic as force can be made
to evaporate --- twice over, in two structurally different ways ---
without changing a single observable prediction is the real content
of this section: it shows "force" to be a feature of a chosen
representation of mechanics, not a fact about the world that every
correct representation must contain.

Hertz's own programme did not survive contact with its critics.
Almost every contemporary reviewer, Boltzmann included, objected
that positing a fresh, permanently unobservable mechanism of hidden
masses for every distinct force law (gravitation, electromagnetism,
\ldots) bought elegance at the price of testability, and the
research programme was essentially abandoned within a generation. It
survives today less as working physics than as an early, and
unusually explicit, demonstration that mechanics has more than one
internally consistent axiomatic foundation --- a lesson later put to
far more productive use.

\subsection{My Own Story}

In the middle of the eighties I studied a problem of quantization
\footnote{Later I describe this problem in more details}.
I have noticed that usual Hamiltonian
\[\mathcal{H}=\Sigma_1^n \frac{\vec p_j^2}{2 m_j} + U(\vec q_j) \quad j=1..n\]
is equivalent to the Hamiltonian defined on expanded phase space 
$(\vec p_j,p_{n+1},\vec q_j,q_{n+1})$:
\begin{equation}
   \label {lnf:hext}
    \mathcal{\tilde H}=\Sigma_1^n \frac{\vec p_j^2}{2 m_j} + U(\vec q_j) \frac{p_{n+1}^2}{2 p_0^2},
\end{equation}
where $p_0$ is arbitrary constant of dimensionality of momentum. 
Indeed, $p_{n+1}$ is integral of motion, fixing its value to be $p_0$ 
we get $ \mathcal{H}=\mathcal{\tilde H}$.

I knew nothing about Hertz's mechanics and called this 'effective mass 
mechanics'. Fortunately due to a number of reasons I did not published
my investigations and so escaped unintentional but clear plagiarism.


\section{The Koopman Formulation of Classical Mechanics}
\label{frm:koop}

\newcommand{\PB}[2]{\{#1,#2\}} 
\newcommand{\CL}{\widehat{\mathcal{L}}}

Classical mechanics is
a theory of nonlinear ordinary differential equations: Hamilton's equations
$\dot q^i = \partial H/\partial p_i$, $\dot p_i = -\partial H/\partial q^i$
generate, in general, a highly nonlinear flow on phase space. It is therefore a
standing surprise, still under-advertised in most mechanics courses, that this
same dynamics admits an exactly equivalent description in which the fundamental
object evolves \emph{linearly}. This is the content of the Koopman formulation:
instead of following the trajectory $x(t)=\phi_t(x_0)$ of a single point
through phase space, one follows the evolution of an entire function $f$ (an
``observable'') under composition with the flow, $f \mapsto f\circ\phi_t$.
Composition with a fixed map is a linear operation on the vector space of
functions, however wild the map $\phi_t$ itself may be. All of the nonlinearity
of the dynamics is absorbed into the definition of a single linear operator, at
the price of moving from a finite-dimensional phase space to an
infinite-dimensional space of functions on it.

This move was introduced by Koopman in 1931 \cite{koopman}, with the
operator-theoretic machinery placed on a rigorous footing almost immediately
afterward by von Neumann \cite{vonNeumann1932} and Stone \cite{stone}, and used
by Koopman and von Neumann together to characterize the different levels of
statistical irregularity (ergodicity, mixing) that a Hamiltonian system can
exhibit, in terms of the spectrum of a single linear operator.
The timing was not an accident: Koopman's paper
is contemporaneous with -- and directly motivated by -- the controversy
surrounding the ergodic hypothesis in statistical mechanics, which was resolved
essentially simultaneously by von Neumann's mean ergodic theorem and Birkhoff's
pointwise ergodic theorem. The operator picture turned the
vague physical question ``does the time average of an observable equal its
phase-space average?'' into a precise question about the point spectrum of a
self-adjoint operator.

\subsection{Phase space, Hamilton's equations, and the flow}

Let phase space be $M=\mathbb R^{2n}$ with coordinates
$x=(q^1,\dots,q^n,p_1,\dots,p_n)$ (everything below goes through, with the
obvious modifications, for a general symplectic manifold $(M,\omega)$). Fix a
Hamiltonian $H:M\to\mathbb R$, assumed smooth and, for simplicity, such that its flow
is complete. Hamilton's equations
\begin{equation}
\dot q^i = \frac{\partial H}{\partial p_i}, 
    \qquad \dot p_i = -\frac{\partial H}{\partial q^i}, \qquad i=1,\dots,n,
\label{hameq}
\end{equation}
define a vector field $X_H$ on $M$ (the Hamiltonian vector field) and hence
a one-parameter group of diffeomorphisms
\begin{equation}
\phi_t : M \to M, \qquad \frac{d}{d t}\phi_t(x) = X_H(\phi_t(x)), 
    \qquad \phi_0=\mathrm{id}, \qquad \phi_t\circ\phi_s=\phi_{t+s}.
\label{flow}
\end{equation}
The Poisson bracket of two functions $f,g:M\to\mathbb R$ is
\begin{equation}
\PB{f}{g} = \sum_{i=1}^n\left(\frac{\partial f}{\partial q^i}
    \frac{\partial g}{\partial p_i}-\frac{\partial f}{\partial p_i}
    \frac{\partial g}{\partial q^i}\right),
\label{PBdef}
\end{equation}
with the sign convention fixed so that, for any smooth $f$, 
its rate of change along a trajectory is
\begin{equation}
\dot f = \PB{f}{H}.
\label{eom}
\end{equation}
(One checks directly from \eqref{PBdef} that 
$\PB{q^i}{H}=\partial H/\partial p_i$ and 
$\PB{p_i}{H}=-\partial H/\partial q^i$, reproducing \eqref{hameq}.)

\begin{proposition}[Liouville's theorem]
\label{liouville}
The Hamiltonian vector field $X_H$ is divergence-free with respect to the 
    Liouville measure $d\mu = \prod_{i=1}^n d q^i\,d p_i$, and 
    consequently $\phi_t$ preserves $\mu$: $\mu(\phi_t(A))=\mu(A)$ 
    for every measurable $A\subset M$.
\end{proposition}
\begin{proof}
$\operatorname{div}X_H = \sum_i\left(\dfrac{\partial}{\partial q^i}
    \dfrac{\partial H}{\partial p_i}+\dfrac{\partial}{\partial p_i}
    \left(-\dfrac{\partial H}{\partial q^i}\right)\right)=
    \sum_i\left(\dfrac{\partial^2 H}{\partial q^i\partial p_i}-
    \dfrac{\partial^2 H}{\partial p_i\partial q^i}\right)=0$ 
    by equality of mixed partials. Measure preservation then follows from 
    Liouville's formula for the Jacobian of the flow of a vector field, 
    $\det D\phi_t(x) = \exp\!\big(\int_0^t 
    \operatorname{div}X_H(\phi_s(x))\,d s\big)=1$.
\end{proof}

This innocuous-looking fact is the load-bearing wall of everything that
follows: it is the only place where the specifically \emph{Hamiltonian}
character of the dynamics enters the Koopman construction, and it is exactly
what will make the Koopman operator unitary rather than merely an isometry or a
general (non-unitary) linear operator.

An \emph{observable} is a (possibly complex-valued) function on phase space, 
$f:M\to\mathbb C$. We work in the Hilbert space
\begin{equation}
\mathcal H = L^2(M,\mu), \qquad \langle f,g\rangle = 
    \int_M \overline{f(x)}\,g(x)\,d\mu(x),
\end{equation}
of observables that are square-integrable with respect to the invariant
Liouville measure. Physically, $L^2(M,\mu)$ is the natural home for observables
when the system carries a normalizable invariant probability distribution
(e.g.\ on a compact energy shell, with $\mu$ the corresponding microcanonical
measure); this is the setting Koopman himself had in mind, motivated by
statistical mechanics. Everything below extends to a general invariant measure
$\mu$, and, with appropriate modifications, to non-invariant or infinite
measures as well.

\begin{definition}[Koopman operator]
For each $t\in\mathbb R$, the Koopman operator $U_t:\mathcal H\to\mathcal H$ 
    is defined by
\begin{equation}
(U_t f)(x) = f(\phi_t(x)), \qquad\text{i.e.}\qquad U_t f = f\circ\phi_t.
\end{equation}
\end{definition}

\begin{proposition}[Linearity and group law]
\label{linear-group}
For all $t,s\in\mathbb R$ and $f,g\in\mathcal H$, $a,b\in\mathbb C$:
\begin{enumerate}
\item $U_t(af+bg) = a\,U_tf+b\,U_tg$;
\item $U_0=\mathrm{id}$ and $U_tU_s=U_{t+s}$; in particular 
    $U_t^{-1}=U_{-t}$.
\end{enumerate}
\end{proposition}
\begin{proof}
(i) is immediate from the definition, since composition with a fixed map 
    $\phi_t$ acts pointwise: $(U_t(af+bg))(x) = af(\phi_t(x))+bg(\phi_t(x))$.
(ii) $(U_tU_sf)(x) = (U_sf)(\phi_t(x)) = f(\phi_s(\phi_t(x))) = f(\phi_{s+t}(x)) 
    = (U_{s+t}f)(x)$, using the group law \eqref{flow} for the flow.
\end{proof}

Statement (i) is the entire point of the construction: $\phi_t$ can be an 
arbitrarily complicated nonlinear diffeomorphism of $M$, yet its pullback 
action $U_t$ on functions is exactly linear, with no approximation involved. 
Statement (ii) says that $\{U_t\}_{t\in\mathbb R}$ is a one-parameter group of linear 
operators on $\mathcal H$.

\begin{proposition}[Unitarity]
\label{unitary}
For every $t\in\mathbb R$, $U_t$ is a unitary operator on $\mathcal H=L^2(M,\mu)$.
\end{proposition}
\begin{proof}
For $f,g\in\mathcal H$,
\begin{equation}
\langle U_tf,U_tg\rangle = 
    \int_M \overline{f(\phi_t(x))}\,g(\phi_t(x))\,d\mu(x).
\end{equation}
Substitute $y=\phi_t(x)$; since $\phi_t$ is a diffeomorphism with 
    $\left|\det D\phi_t\right|=1$ by Proposition~\ref{liouville}, 
    the measure transforms as $d\mu(x)=d\mu(y)$, so
\begin{equation}
\langle U_tf,U_tg\rangle = \int_M \overline{f(y)}\,g(y)\,d\mu(y) =
    \langle f,g\rangle.
\end{equation}
Thus $U_t$ is an isometry; since it is also invertible 
    (Proposition~\ref{linear-group}(ii), with inverse $U_{-t}$, 
    itself an isometry by the same argument), $U_t$ is unitary.
\end{proof}

It is worth being explicit about what did the work: unitarity is \emph{not} an
automatic feature of ``Koopman operators'' for arbitrary dynamical systems; it
requires an invariant measure. For a general (e.g.\ dissipative) dynamical
system with no invariant probability measure, the Koopman operator built the
same way is still linear and still a composition operator, but it need not be
unitary or even bounded on a naturally chosen $L^2$ space. It is Liouville's
theorem, i.e.\ the specifically Hamiltonian structure of classical mechanics,
that upgrades ``linear'' to ``unitary.''

Finally, strong continuity ($t\mapsto U_tf$ continuous in $\mathcal H$ for each
fixed $f$, under mild regularity of $H$ and $f$) makes $\{U_t\}$ a
\emph{strongly continuous one-parameter unitary group}, which is precisely the
hypothesis of Stone's theorem, taken up next.

\begin{definition}[Generator]
The (infinitesimal) generator of the Koopman group is the operator
\begin{equation}
Df := \left.\frac{d}{d t}\right|_{t=0} U_tf,
\end{equation}
defined on the (dense) domain of $f\in\mathcal H$ for which the limit 
    exists in $\mathcal H$.
\end{definition}

By the chain rule and \eqref{eom},
\begin{equation}
(Df)(x) = \left.\frac{d}{d t}\right|_{t=0} 
    f(\phi_t(x)) = X_H f(x) = \PB{f}{H}(x),
\end{equation}
so the generator is simply the Poisson bracket with $H$:
\begin{equation}
D = \PB{\,\cdot\,}{H}.
\label{Ddef}
\end{equation}
Because $\phi_t$ is a flow, $D$ commutes with $U_t$ and one has, 
formally,
\begin{equation}
U_t = e^{tD}, \qquad \frac{d}{d t}U_t = D\,U_t = U_t\,D.
\label{Uexp}
\end{equation}

\begin{proposition}[Skew-adjointness of $D$]
\label{skew}
On functions vanishing at infinity (or, if $M$ is compact, with no further 
    hypotheses), $D$ is skew-adjoint with respect to 
    $\langle\,\cdot\,,\,\cdot\,\rangle$: $\langle Df,g\rangle =
    -\langle f,Dg\rangle$ for all $f,g$ in the domain.
\end{proposition}
\begin{proof}
Take $f,g$ real-valued and smooth with compact support, for definiteness 
    (the general complex case follows by linearity, using that $D$ has real
    coefficients so $D\bar f=\overline{Df}$). Integrating by parts twice in 
    \eqref{PBdef} (once in $q^i$, once in $p_i$, discarding boundary terms) 
    gives
\begin{align}
\int_M \PB{f}{H}\,g\,d\mu
&= \sum_i\int_M\left(\frac{\partial f}{\partial q^i}\frac{\partial H}
    {\partial p_i}-\frac{\partial f}{\partial p_i}\frac{\partial H}
    {\partial q^i}\right)g\,d\mu \notag\\
&= -\sum_i\int_M f\left(\frac{\partial^2H}{\partial q^i\partial p_i}g
    +\frac{\partial H}{\partial p_i}\frac{\partial g}{\partial q^i}
    -\frac{\partial^2H}{\partial p_i\partial q^i}g
    -\frac{\partial H}{\partial q^i}\frac{\partial g}
    {\partial p_i}\right)d\mu.
\end{align}
The two mixed-partial terms cancel identically, leaving
\begin{equation}
\int_M \PB{f}{H}\,g\,d\mu = -\int_M f\left(\sum_i\Big(\frac{\partial H}
    {\partial p_i}\frac{\partial g}{\partial q^i}-
    \frac{\partial H}{\partial q^i}\frac{\partial g}
    {\partial p_i}\Big)\right)d\mu = \int_M f\,\PB{g}{H}\,d\mu,
\end{equation}
i.e.\ $\langle Df,g\rangle = -\langle f, Dg\rangle$ (the last relative sign is
    exactly $-1$ because $\PB{g}{H}=-\PB{H}{g}$ is being compared to 
    $D g = \PB{g}{H}$; the displayed identity already contains it).
\end{proof}

A skew-adjoint operator becomes self-adjoint (Hermitian) after multiplication 
by $i$. Define the \emph{Liouvillian}
\begin{equation}
\boxed{\CL := -i D = -i\,\PB{\,\cdot\,}{H} = i\,\PB{H}{\,\cdot\,}.\ }
\label{Liouvillian}
\end{equation}
By Proposition~\ref{skew}, $\CL^\dagger=(-i D)^\dagger = i D^\dagger = 
i(-D)=-i D=\CL$: the Liouvillian is Hermitian. Combined with \eqref{Uexp}, 
and since $D=i\CL$,
\begin{equation}
U_t = e^{tD} = e^{i t\CL}.
\label{Ut-CL}
\end{equation}
This is precisely Stone's theorem \cite{stone} at work: a strongly continuous
one-parameter \emph{unitary} group is always of the form $e^{i t A}$ for a
unique self-adjoint (generally unbounded) operator $A$; here $A=\CL$, built
entirely out of the classical Poisson bracket with $H$. Equation \eqref{Ut-CL}
is the precise sense in which classical Hamiltonian evolution is ``already''
Schr\"odinger-like, well before any quantization has taken place. Note,
however, that conventions for the sign and placement of $i$ in $\CL$ vary
across the literature; \eqref{Liouvillian} is fixed here purely so that
\eqref{Ut-CL} and the density equation of the next section come out in the most
familiar-looking form.


The Koopman operator $U_t$ propagates \emph{observables} forward in time: it is
the classical analogue of the Heisenberg picture. There is a dual,
Schr\"odinger-like picture that propagates \emph{probability densities},
obtained simply by taking the adjoint.

\begin{definition}[Transfer, or Perron--Frobenius, operator]
The operator $P_t := U_t^\dagger$ acting on densities 
    $\rho\in\mathcal H$ is defined by duality,
\begin{equation}
\int_M (P_t\rho)\,f\,d\mu = \int_M \rho\,(U_tf)\,d\mu 
    \qquad \text{for all } f\in\mathcal H.
\end{equation}
\end{definition}

Since $U_t$ is unitary, $P_t = U_t^\dagger = U_t^{-1}=U_{-t}$, so $(P_t\rho)(x)
= \rho(\phi_{-t}(x))$: the density at $x$ at time $t$ equals the density that
was at the point $\phi_{-t}(x)$ that flowed \emph{into} $x$. This is exactly
the statement that probability is transported (not created or destroyed) along
the flow, consistent with Proposition~\ref{liouville}.

Differentiating $P_t\rho = \rho\circ\phi_{-t}$ at $t=0$, the generator of $P_t$ 
is $-D$, i.e.
\begin{equation}
\frac{\partial \rho_t}{\partial t} = -D\rho_t = -\PB{\rho_t}{H} = 
    \PB{H}{\rho_t},
\label{liouville-eq}
\end{equation}
which is exactly the classical Liouville equation, usually written 
$\partial_t\rho+\PB{\rho}{H}=0$. Multiplying \eqref{liouville-eq} by $i$ and 
using $D=i\CL$,
\begin{equation}
\boxed{\ i\,\frac{\partial \rho_t}{\partial t} = \CL\,\rho_t. \ }
\label{kvn-density}
\end{equation}
This is a genuine Schr\"odinger-type equation for the classical phase-space
density, with Hermitian ``Hamiltonian'' $\CL$: $\rho_t = e^{-i t \CL}\rho_0 =
P_t\rho_0$. It is the direct forerunner of the Koopman--von Neumann equation. 

\begin{remark}
The two pictures are related exactly as in quantum mechanics: observables
    evolve by $U_t=e^{i t\CL}$ (formally, $i\partial_tf_t = -\CL f_t$, the
    classical Heisenberg equation, consistent with $\dot f=\PB{f}{H}$), while
    densities evolve by the inverse/adjoint operator $P_t=e^{-i t\CL}$,
    Eq.~\eqref{kvn-density}. The relative sign is the same
    Heisenberg/Schr\"odinger duality familiar from quantum theory, here derived
    from nothing more than the unitarity of $U_t$.
\end{remark}

The Koopman formulation adds no new physics to classical mechanics: every
statement derivable in the operator language is, by construction, a restatement
of Hamilton's equations. What it buys is a change of mathematical category,
from nonlinear dynamics on a finite-dimensional manifold to linear operator
theory on an infinite-dimensional Hilbert space, and that change of category is
not free -- it is purchased with the invariant measure furnished by Liouville's
theorem, which is what makes the Koopman operator unitary rather than merely
linear. Once purchased, the machinery of spectral theory becomes directly
applicable to a question -- how ergodic, or how mixing, is this system? -- that
has no simple formulation in the original nonlinear language, and this is
precisely the use Koopman, von Neumann, and Birkhoff made of it in the early
1930s.
 
